% =====================================================================
%  Conferenze, GAMM Roma (Checkpoint Festival, 3a edizione) — Panel 3
%  Dieci anni di horror indipendente: dal fan remake a Capcom a uno
%  studio di provincia
%  Fonte: trascrizione automatica "Track 3". Intervista-panel con i tre
%  fondatori di Invader Studios (Michele, Alessandro, Tiziano), condotta
%  da una moderatrice. La registrazione comincia a domanda già avviata;
%  durante la proiezione dei trailer la trascrizione riporta ripetizioni
%  meccaniche dell'audio dei video, che non si riportano.
%  Etichette con prefisso "c1-3:". Nessun pacchetto aggiuntivo.
% =====================================================================

\chapter[Dieci anni di horror indipendente]{Dieci anni di horror
indipendente: da un fan remake a uno studio di provincia}
\label{cap:c1-3-horror-indipendente}

Il terzo incontro al GAMM è un'intervista ai tre fondatori di uno studio
italiano specializzato in \textbf{survival horror}, \textbf{Invader
Studios}\footnote{Il nome dello studio è trascritto ``Independent Studios'';
l'identificazione con Invader Studios è sicura dal contesto (\emph{Daymare:
1998}, trascritto ``Te nel 1998'', ``Damer'', ``Danger''; \emph{Daymare: 1994
Sandcastle}, trascritto ``Sencasso'', ``Seikastra'', ``Dead Castle'',
``Damer94''; la sede in provincia di Roma, a Olevano Romano, trascritta
``Dolevano Romano'').}, nato come gruppo di appassionati di \emph{Resident
Evil 2}. La moderatrice ripercorre con loro la storia dello studio, le scelte di
design dei loro giochi, l'esperimento di un titolo cooperativo dai toni comici e
gli obiettivi per i prossimi dieci anni.

I tre ospiti, chiamati per nome di battesimo:
\begin{description}
  \item[Michele] (Michele Giannone\footnote{Il cognome non è pronunciato nella
        parte registrata; si tratta del CEO e cofondatore dello studio. Altrove
        la trascrizione lo chiama anche ``Miguel''.}), alla guida
        dell'azienda;
  \item[Alessandro], game designer, che cura anche gli aspetti di design legati
        al suono;
  \item[Tiziano], artista 3D, oggi soprattutto sulle ambientazioni.
\end{description}

Struttura del capitolo:
\begin{enumerate}
  \item da \emph{Resident Evil 2} a Osaka (sezione~\ref{sec:c1-3-capcom});
  \item il primo incontro con \emph{Resident Evil 2}
        (sezione~\ref{sec:c1-3-imprinting});
  \item la paura sulla soglia di una stanza (sezione~\ref{sec:c1-3-soglia});
  \item quando la tensione non funziona: iterazione e messaggi ambientali
        (sezione~\ref{sec:c1-3-iterazione});
  \item riscrivere le regole: l'horror che fa ridere
        (sezione~\ref{sec:c1-3-panic});
  \item il marchio di fabbrica e le aspettative del pubblico
        (sezioni~\ref{sec:c1-3-marchio} e~\ref{sec:c1-3-aspettative});
  \item i prossimi dieci anni e la storia per immagini
        (sezioni~\ref{sec:c1-3-futuro} e~\ref{sec:c1-3-foto}).
\end{enumerate}

% ---------------------------------------------------------------------
\section{Da Resident Evil 2 a Osaka}
\label{sec:c1-3-capcom}

La moderatrice riassume l'origine: un gruppo di appassionati di
\emph{Resident Evil}, e in particolare di \emph{Resident Evil 2}, comincia a
lavorare a un progetto amatoriale (un \emph{fan remake}\footnote{Il progetto non
è nominato esplicitamente nella parte registrata; si tratta del remake
amatoriale di \emph{Resident Evil 2} su cui il gruppo lavorava prima della
fondazione dello studio.}). Il progetto viene \textbf{bloccato} perché il gruppo
è chiamato a \textbf{Osaka}, ``da niente poco di meno che \textbf{Capcom}''. E
dal ritorno da Osaka comincia la vera storia dello studio.

\subsection{L'incontro con Capcom}

Il gruppo sapeva che altri progetti amatoriali erano stati fermati \textbf{per
vie legali}, e da un'azienda giapponese si aspettava qualcosa di simile. Capcom
invece fu ``molto più attenta''. Il momento era delicato per la saga: secondo
Michele \emph{Resident Evil} stava ``prendendo una piega un po' verso il
basso'', \emph{Resident Evil 7} non era ancora stato annunciato, né tantomeno
il remake; Capcom era molto attenta a dialogare con la community e ``cavalcò''
il piccolo successo del gruppo.

L'approccio non fu una diffida, ma un \textbf{apprezzamento del lavoro}. Alcune
idee proposte dal gruppo, racconta Michele, sono poi \textbf{confluite nel titolo
ufficiale}, e i membri del gruppo sono stati inseriti nei \textbf{ringraziamenti
ufficiali del gioco}\footnote{La trascrizione dice ``nei premi su ufficiali di
gioco'': con ogni probabilità ``nei credits ufficiali'' o ``nei ringraziamenti
speciali''.}.

A Osaka si sedettero con alcuni producer della serie\footnote{I nomi sono
storpiati (``Giunta che Luigi, Kirabayashi\dots''); tra questi è riconoscibile
Yoshiaki Hirabayashi, producer del remake di \emph{Resident Evil 2}.} e
mostrarono la loro visione di un \emph{Resident Evil} basata sulle proprie idee,
con video e documentazione. La risposta: ``le idee sono belle, ma ovviamente non
possiamo darvi un contributo di nessun tipo''. Ci fu delusione (``eravamo
giovani, non potevi aspettarti chissà quanto di più''), ma anche uno
\textbf{sprone}: ``tornate a casa e fate qualcosa di vostro con quello che ci
avete mostrato, non legato a \emph{Resident Evil}''. Erano quattro o cinque
persone, senza esperienza nel creare da zero un'esperienza completa, ma con
studi ed esperienze in altri studi. Quella fu ``la scintilla'': il riconoscimento
che le loro idee avevano validità.

Il lavoro successivo fu \textbf{trasformare la passione in proprietà
intellettuale}: sradicare l'idea da ogni legame con \emph{Resident Evil} e farne
qualcosa di loro. Il risultato fu \emph{Daymare: 1998}.

% ---------------------------------------------------------------------
\section{L'imprinting: la prima volta con Resident Evil 2}
\label{sec:c1-3-imprinting}

Domanda ``da videogiocatrice'': vi ricordate la prima volta che avete giocato
\emph{Resident Evil 2}, nel 1998?

\paragraph{Alessandro.} È stato ``uno dei pochi imprinting fondamentali'' della
sua carriera. La sua storia di giocatore era cominciata prima, con le console
(Commodore 64, Nintendo) e le \textbf{sale giochi}, oggi quasi scomparse (il
museo, osserva, prova a recuperare quella connessione). Ma con \emph{Resident
Evil 2} cambiò ``la linea'': si passava da un intrattenimento molto ludico a
\textbf{videogiochi più maturi}, con temi come la vita, la morte, l'amore,
emozioni che un ragazzino difficilmente trovava nei giochi arcade. ``Quello fu
proprio un pugno'': da lì si è formato tutto, ``e non a caso'' quindici anni dopo
hanno deciso di fare un gioco intitolato al 1998.

\paragraph{Michele.} ``Faccio il vecchio della situazione'': c'erano le
\textbf{demo sulle riviste}. Ricorda l'immagine usata moltissimo nel marketing
del gioco (una mano, un volto) e la demo, con un \textbf{conto alla rovescia}:
si poteva giocare solo dieci minuti, poi il tempo scadeva. Per lui fu ``un pugno
in faccia'': riusciva finalmente a creare una connessione tra i film di
\textbf{George Romero} sugli zombie e qualcosa che poteva giocare. Veniva da
giochi colorati e divertenti di un'altra tipologia; ``l'horror è entrato e non è
più uscito''.

\paragraph{Tiziano.} Anche lui, come Michele, era appassionato dei film di zombie
di Romero, visti ``di traverso'' a un'età in cui non avrebbe dovuto, e che lo
spaventarono tanto da creare la passione. \emph{Resident Evil} era la
congiunzione tra videogiochi e zombie. Ricorda soprattutto l'\textbf{impatto
grafico}: i \textbf{fondali pre-renderizzati}, già usati in precedenza, ma con
il secondo capitolo ``un balzo avanti importante''. Ambientazioni che oggi non
sembrerebbero nitide, ma che all'epoca erano ``qualcosa di incredibile'', con
personaggi messi in scena in quelle inquadrature particolari.

\subsection{Il primo personaggio di Tiziano}

La moderatrice chiede a Tiziano, ``prima ancora che sviluppatore, artista'',
quale sia stato il primo personaggio che ha creato e di cui è ancora orgoglioso.
Tiziano confessa di non essere ``tendenzialmente orgoglioso'' di ciò che fa; i
primi lavori li ha tolti da internet. Erano modelli legati al \textbf{Signore
degli Anelli} (armature, un \textbf{troll delle caverne} ``non uscito come quello
del film''): si divideva tra zombie e fantasy. Era combattuto tra fare
\textbf{personaggi} o \textbf{ambientazioni}, e ha provato entrambe le cose; poi,
quando ha capito che l'industria cercava molto più \textbf{environment artist}
che \emph{character artist}, ha preso quella strada. La passione per i
personaggi è rimasta, e ogni tanto in azienda ``qualche task sui personaggi me
la prendo'', anche se ormai fanno fare quasi tutto ai ragazzi del team.

% ---------------------------------------------------------------------
\section{Sulla soglia di una stanza}
\label{sec:c1-3-soglia}

La moderatrice propone ad Alessandro un esercizio: ``sono una videogiocatrice,
entro in una stanza di \emph{Sandcastle} e mi fermo sulla soglia. Come agisci,
con il sound design e con il game design, per farmi percepire paura e
attenzione anche se sono immobile, e al tempo stesso per invogliarmi a
entrare?''

Alessandro premette una regola che considera tra le principali del game design:
\begin{quote}
  stabilire delle \textbf{regole molto chiare} e dare al giocatore la
  possibilità di averle chiare mentre gioca.
\end{quote}
C'è una prima fase in cui il giocatore le comprende; possono anche cambiare nel
corso dell'avventura, ma sono \textbf{i pilastri} su cui si regge l'intero gioco.

In \emph{Sandcastle} la regola è precisa: \textbf{non ci sono nemici già
posizionati in scena}; possono arrivare \textbf{in qualsiasi momento}, grazie a
un sistema di \textbf{teletrasporto}. È uno spunto preso da \textbf{Dead Space},
dove le minacce non erano leggibili né presenti in scena. Quindi, sulla soglia,
il giocatore \textbf{sa} che può essere attaccato in ogni momento; deve però
proseguire per finire l'avventura, attraversando zone in cui non riesce a
percepire il pericolo finché non se lo trova davanti.

È un horror diverso da quello \textbf{ansiogeno}, che mostra prima ciò che ti
aspetta e ti fa combattere con le tue paure prima dell'azione. Con
\emph{Sandcastle} hanno scelto un horror più \textbf{adrenalinico},
\textbf{più action}: un \textbf{picco} di paura e adrenalina nel momento esatto
in cui entri e ti imbatti negli avversari.

% ---------------------------------------------------------------------
\section{Quando la tensione non funziona}
\label{sec:c1-3-iterazione}

Domanda a tutti e tre: c'è stato un momento in cui gli elementi (suono, design,
arte) non si coordinavano e non riuscivate a creare la tensione voluta? Come
l'avete risolto?

\paragraph{``È sempre così'' (Michele).} Quando si mette a terra il
\emph{design document} lo si fa con molte aspettative e con il bagaglio dei
progetti precedenti: si ha una \textbf{previsualizzazione} di ciò che si vuole.
Poi subentrano tutti gli aspetti tecnici e artistici che contribuiscono alla
messa in scena, e servono moltissime \textbf{revisioni}. La potenza del
videogioco è che è \textbf{un medium non lineare}: ognuno lo gioca a modo suo.
L'obiettivo è riuscire \textbf{statisticamente} a far vivere a più giocatori
possibili lo stesso momento nello stesso modo, e questo passa da un
\textbf{processo iterativo}: ``questa cosa non funziona, la rifacciamo, la
rifacciamo, la rifacciamo'', finché statisticamente l'effetto si riproduce. Nei
giochi \textbf{lineari} è più facile; in quelli \textbf{più aperti} è molto più
difficile portare il giocatore nel punto che vuoi, nel momento in cui vuoi.

Con ogni gioco lo studio diventa più ambizioso, quindi le difficoltà aumentano.
\emph{Daymare: 1998} era il titolo che ``secondo tutti non era possibile fare con
così poche persone'': ci hanno sbattuto la testa, hanno cambiato scelte e
imparato moltissimo. È stata ``una scuola'', a metà tra incoscienza, passione,
volontà e conoscenze. Oggi l'esperienza aiuta a prendere decisioni migliori, e il
feedback del pubblico alimenta i progetti successivi; ma ``non si arriva mai al
punto'': è importante non smettere di migliorare.

\paragraph{I messaggi affidati alle ambientazioni (Tiziano).} Lo studio ha una
serie di regole per comunicare ciò che vuole (horror, \emph{jump scare}, momenti
riflessivi), ma non sempre ciò che si ha in mente si traduce come previsto.
Lavorando con Michele sulle ambientazioni, cercano di \textbf{mandare messaggi
con l'ambiente}; a volte però i ragionamenti sono ``troppo machiavellici'',
caricano l'ambiente di messaggi troppo complessi da comunicare con la sola
disposizione di pochi elementi di scena che raccontano che cosa è successo prima
dell'arrivo del protagonista.

Cercano di \textbf{contestualizzare tutto} e di non lasciare nulla al caso. Un
esempio nato proprio da \emph{Resident Evil 2}: nella stazione di polizia del
gioco \textbf{non hanno mai trovato i bagni}. Da allora, nei loro giochi, ``se
ci sono dodici stanze per le conferenze ci devono stare altrettanti bagni'': fanno
uno studio su \textbf{quanti bagni servono per quante persone} dovrebbero
frequentare quell'ambiente. Poi si accorgono che a volte questo va a scapito
dell'aspetto ludico, e devono riconcentrarsi su \textbf{come il giocatore gioca}
piuttosto che sul numero di bagni. ``È divertente, lo facciamo spesso'', ma non
sempre una sottotrama affidata all'ambiente viene colta.

% ---------------------------------------------------------------------
\section{Riscrivere le regole: l'horror che fa ridere}
\label{sec:c1-3-panic}

Tutti hanno parlato di regole. Ma a un certo punto lo studio le ha
\textbf{riscritte}, passando a un titolo cooperativo dai toni comici\footnote{Il
titolo è trascritto ``Panic Demonic'' e poi solo ``Panic''; non è stato possibile
verificarne la grafia esatta.}. Come?

\subsection{Perché farlo}

Michele indica più ragioni:
\begin{itemize}
  \item il team aveva bisogno di \textbf{una pausa} dai giochi più seriosi, ed
        era giusto provare qualcosa di più scanzonato;
  \item ragioni \textbf{commerciali}: lo studio voleva provare
        l'\textbf{autopubblicazione}, che su un titolo di punta con grandi
        investimenti è molto rischiosa per una realtà come la loro; su un titolo
        più sperimentale si mettono in campo risorse diverse e ``te lo puoi
        permettere''.
\end{itemize}
Hanno scoperto che \textbf{l'horror può anche essere divertente}: la sfida era
fare cose che mettono paura \emph{mentre} fanno ridere con gli amici, il filone
dei giochi cooperativi horror nato su Steam. Anche il lavoro sul suono è
cambiato: i suoni pensati per spaventare non funzionavano, perché ora dovevano
anche divertire. Il cambio di prospettiva ha dato idee utili anche per i progetti
più seriosi.

Per Tiziano è stato anche un segno dell'\textbf{evoluzione del posizionamento
aziendale}:
\begin{itemize}
  \item \emph{Daymare: 1998}: hardcore, con limitazioni tecniche ma anche di
        meccaniche;
  \item \emph{Daymare: 1994 Sandcastle}: più aperto all'action, nemici e
        ambientazioni diversi; ancora un \emph{action horror} in terza persona;
  \item il nuovo titolo: \textbf{prima persona}, completamente aperto a un
        pubblico e a interazioni diverse, e il primo \textbf{autopubblicato}
        (prima avevano sempre lavorato con publisher).
\end{itemize}
È la dimostrazione, a se stessi e ai diversi pubblici, di saper fare cose
diverse.

\subsection{Scelte contrastanti con l'horror}

Alessandro sottolinea che in questo gioco hanno fatto scelte di design e di arte
\textbf{in contrasto con il tema horror}:
\begin{itemize}
  \item la grafica è volutamente \textbf{cartoonish};
  \item il gioco \textbf{in cooperativa} è il contrario di come si costruisce la
        tensione: \textbf{più si è, meno l'horror ha effetto}.
\end{itemize}
Tuttavia c'è un momento chiave: puoi \textbf{restare da solo}. E allora
``improvvisamente ti arriva tutta quella carica horror'' a cui i giocatori non
erano abituati. ``Divertitevi, scherzate, parlatevi: ma il gioco, a un certo
punto, vi metterà nella condizione di avere improvvisamente paura.''

\subsection{Gestire il caos dei giocatori}

Come avete trovato la formula, in singolo e in gruppo, quando i giocatori possono
anche autosabotarsi? Michele: su questo titolo hanno deciso di
\textbf{togliere} alcune regole. Di solito cercano di non permettere al giocatore
di fare troppo; qui invece \textbf{lo incentivano}. In fase di design hanno
pensato molto più al multigiocatore che al singolo, che si sperimenta comunque
quando ci si ritrova soli.

Però non volevano un caos totale, perché il caos toglie attenzione alle situazioni
che il gioco prepara: se stai scherzando e succede qualcosa che non ti aspetti,
arriva il \emph{jump scare}. L'esempio è un mostro che non tutti hanno visto:
quando un giocatore si allontana e muore, il mostro \textbf{prende le sue
sembianze}; tu lo incontri, pensi che sia il tuo compagno, e quando te ne accorgi
``è troppo tardi''.

\subsection{La critica sociale}

Tiziano aggiunge il concept: il genere umano è finito e c'è una nuova società
governata dai \textbf{mostri}, che aprono un'azienda di \textbf{spedizioni}; il
giocatore è un \textbf{postino umano}, sfruttato e sacrificabile. È una
\textbf{critica sociale} alla sacrificabilità di certe categorie di lavoratori.
Il concetto è già caotico in sé (la \textbf{ribellione del lavoratore}), ma a un
certo punto il ``padrone'' si mostra più forte e subentrano la paura e l'horror:
``devi stare attento anche a come ti ribelli''.

Durante il panel viene proiettato il trailer del gioco\footnote{La trascrizione
riporta più volte di seguito una battuta del trailer (sulla ``blue map'' che
permette di trovare le stanze segrete senza sprecare bombe): è l'audio del video,
ripetuto meccanicamente. Vengono poi mostrati i trailer di lancio di
\emph{Daymare: 1998}, uscito su PC nel settembre 2019 e su console, in Giappone,
nel febbraio 2020, con un'edizione fisica giapponese di 1.000 copie, secondo
quanto dice il relatore.}. Michele ricorda con orgoglio che \emph{Daymare: 1998}
è stato fatto ``praticamente in cinque persone'', il titolo che ``tutti dicevano
impossibile''; ancora oggi ``fa bene''.

% ---------------------------------------------------------------------
\section{Il marchio di fabbrica}
\label{sec:c1-3-marchio}

Qual è l'elemento che il pubblico riconosce subito come vostro?

\paragraph{Michele.} Da ciò che scrivono i giocatori per mail e sui social:
riuscire a essere \textbf{in parte derivativi} (da \emph{Resident Evil} e da
tutto ciò che li ha segnati) ma \textbf{originali} allo stesso tempo, creando un
contesto e un mondo che funzionano al di là dei rimandi. Il genere funziona così:
se un gioco o un film ti piace, cerchi subito qualcosa di simile, ``e i numeri
ultimamente ce lo dicono''. Lo studio è passato da \textbf{5 a quasi 25
persone} e oggi lavora internamente a ogni aspetto: ambientazioni, narrativa,
testi, dialoghi, \emph{cutscene}, animazione, \textbf{audio} (fondamentale e
complessissimo in un horror). C'è un clima ``come una famiglia'': le persone sono
contente di venire a lavorare, e questo si sente in un prodotto che, pur con i
suoi difetti, ha un'impronta non scontata.

\paragraph{Tiziano.} Vorrebbe che lo studio fosse associato a un
\textbf{revamp degli anni Ottanta--Novanta}. Oggi è scontato (si pensi a
\emph{Stranger Things}), ma quando hanno cominciato, nel 2013--2014, riprendere
qualcosa di considerato ``vecchio'' non lo era: la stessa Capcom aveva messo da
parte i suoi \emph{Resident Evil} storici. La chiave è la \textbf{nostalgia}:
quella che vivono loro e che vogliono trasmettere.

\paragraph{Alessandro.} La missione iniziale era proprio \textbf{riproporre
esperienze, sapori e sensazioni} trovate in certi titoli e momenti della loro
vita. I commenti più belli sono di chi ritrova \textbf{emozioni dimenticate}. Non
è solo l'aspetto visivo, ma il \textbf{modo di narrare}, di mettere in scena e di
veicolare i messaggi: anche un'implementazione moderna, per la cura di certi
dettagli, ti fa percepire il gioco come ``erede'' di quei classici.

% ---------------------------------------------------------------------
\section{Le aspettative del pubblico}
\label{sec:c1-3-aspettative}

Vi siete mai sentiti limitati dalle aspettative del pubblico, che si aspetta da
voi un certo genere?

Michele ammette che il titolo cooperativo è stato \textbf{divisivo} per la
community, ma era un esperimento che aveva senso fare: ``se un'idea secondo noi è
valida, la mettiamo a terra''. Il vero limite viene dall'\textbf{ambizione}:
quando si punta in alto, si riescono a \textbf{mascherare} i limiti di una
produzione a basso budget, e allora il pubblico pensa che tu sia un doppia A o
un tripla A, ``e non è così''. Si aspetta valori produttivi che lo studio non può
ancora permettersi. La tecnologia degli ultimi anni aiuta, lo studio cresce in
\emph{know-how} e mostra sempre di più; e chi capisce la filosofia dello studio
diventa un fan.

Michele racconta poi l'esercizio che ``tutti i creativi fanno'': prendere le
critiche, mettere da parte quelle positive (``non servono a nulla'') e
\textbf{analizzare a fondo quelle negative o costruttive}. Il giocatore ha il
diritto di giudicare senza sapere chi c'è dietro; quando invece fa lo sforzo di
capirlo, diventa un fan, perché comprende che non ci sono ``i mezzi milionari''
delle grandi case.

Lo studio è \textbf{di provincia}, ``orgogliosamente'': viene da un territorio che
produce olio e vino, dove quando dicono che lavoro fanno ``non comprendono
minimamente''. Per loro è un vanto poter mostrare che da quell'ecosistema si può
arrivare a entrare in Capcom a discutere di \emph{Resident Evil}. Il messaggio
che vogliono comunicare è chiaro:
\begin{quote}
  un gioco che assomiglia a \emph{Resident Evil} non ha i valori produttivi di
  \emph{Resident Evil}. Se accettate questo compromesso, vi potete divertire.
\end{quote}
E i giocatori lo accettano \textbf{se il prezzo è coerente}: ``se costi come
\emph{Resident Evil}, è giusto che ti critichino come \emph{Resident Evil}; se
costi un terzo, è giusto che le critiche siano misurate sulla proposta''. La
community, molto attenta, a volte si accende in discussioni tra fan nelle
recensioni; il seguito, dice, è più internazionale che italiano.

% ---------------------------------------------------------------------
\section{I prossimi dieci anni}
\label{sec:c1-3-futuro}

Lo studio ha ormai dieci anni (``16 giugno'', precisa uno dei tre). Quali
obiettivi per i prossimi dieci?

\paragraph{Un nuovo titolo con un partner importante.} Il primo obiettivo è già
spuntato: un nuovo gioco, non ancora annunciato, con un partner internazionale.
Michele spiega che, quando si parla di progetti futuri, prima che arrivi un
contratto firmato passano anche \textbf{7--8 mesi}, con ``tanti buchi
nell'acqua''; questa volta invece si è arrivati alla firma e all'inizio dello
sviluppo. Il titolo li posizionerà in modo importante a livello internazionale,
``un sogno che si avvera''; ne hanno altri da esplorare, con canali che sembrano
aprirsi\footnote{In questo tratto la trascrizione è quasi illeggibile, con
lunghe ripetizioni meccaniche e frammenti (si accenna a un progetto ``basato
su'' qualcosa e a un ``attore molto importante''). Le informazioni sul nuovo
progetto sono comunque riservate e i relatori evitano di darne.}.

\paragraph{Le persone e il luogo (Alessandro).} In dieci anni vorrebbe pensare
allo studio come in realtà già è: un luogo in cui molte persone hanno scelto di
\textbf{trasferirsi in provincia}, facendo il percorso inverso a quello di tanti
ragazzi costretti a lasciare i paesi per studiare o lavorare altrove. In un
periodo di \textbf{profonda trasformazione} dell'industria, con i licenziamenti
(i \emph{layoff}), l'obiettivo è \textbf{consolidare} la struttura: dare a chi
lavora con loro progetti stimolanti ma anche \textbf{certezza}, che in questa
industria ``non è mai data''.

\paragraph{Non diventare troppo grandi (Michele).} Paradossalmente bisogna
\textbf{non crescere troppo}. Crescere sì, coinvolgere più persone, ma con i passi
giusti: collaborando con altre aziende hanno visto che più sei grande, più
problemi nascondi e meno ci pensi. ``Dieci anni fa, nelle nostre più rosee
aspettative, ci vedevamo a cento persone; da 25 magari saremo 50, e ci
accontenteremo'': fare il lavoro che piace, che permette di vivere, con progetti
``delle giuste dimensioni''.

% ---------------------------------------------------------------------
\section{La storia dello studio per immagini}
\label{sec:c1-3-foto}

Il panel si chiude con una serie di foto e video commentati:
\begin{itemize}
  \item il \textbf{primo studio} (2015), un garage con postazioni ``prese dalle
        rispettive camerette'' da varie regioni, spesso allagato, con un solo
        bagno, dove si lavorava ``giorno e notte'';
  \item il \textbf{secondo studio}, un appartamento (``eravamo migliorati: due
        bagni''), dove sono cresciuti e dove i primi ragazzi si sono
        trasferiti;
  \item le prime fiere (Milan Games Week, Gamescom nel 2018\footnote{La
        trascrizione dice ``James Week, la Chiesa con il 2018'': la lettura è
        un'ipotesi.}) e la ricerca di un publisher;
  \item il primo viaggio in Giappone per \emph{Daymare: 1998}, il nuovo incontro
        con Capcom e con \textbf{Shinji Mikami}\footnote{Trascritto ``Cendrick
        Mikami''.}, che definiscono ``il nostro mentore'';
  \item una foto ``surreale'': i \textbf{negozi giapponesi allestiti} con il
        materiale di \emph{Daymare} durante il lancio promozionale, di cui non
        avevano capito l'impatto;
  \item poco prima del Covid, un \textbf{press tour giapponese} organizzato dal
        distributore del gioco in Giappone, che portò nello studio, in provincia,
        giornalisti di Kotaku, IGN Japan e Famitsu: ``come se avessimo aperto
        casa'', con un approccio caloroso e ``molto italiano'', ricordato ancora
        oggi;
  \item un incontro a casa di un autore legato a \emph{Dead Space}, che li ha
        accolti con una bistecca e ha raccontato la sua idea di un
        \emph{Dead Space 4} mai realizzato\footnote{Il passaggio è molto
        confuso (``Hertz Space'', ``Slovenia con la bistecca''); si riporta solo
        il senso generale, senza poter identificare la persona.};
  \item il Tokyo Game Show più recente e le copie del gioco nei negozi
        giapponesi nel 2024: ogni volta che vanno in Giappone comprano tutti i
        giochi usati che trovano;
  \item infine il video dello \textbf{studio attuale}, ``ancora vuoto'' al
        momento del trasloco e oggi pieno, con tutti i reparti interni
        (programmazione, animazione, \emph{motion capture}, audio).
\end{itemize}

% ---------------------------------------------------------------------
\section*{Punti chiave del capitolo}
\addcontentsline{toc}{section}{Punti chiave del capitolo}

\begin{itemize}
  \item Un fan remake di \emph{Resident Evil 2} viene fermato, ma Capcom invita
        il gruppo a Osaka e lo sprona a creare \textbf{una propria IP}: nasce
        \emph{Daymare: 1998}, fatto in cinque persone.
  \item Regola cardine del design: \textbf{regole chiare} che il giocatore
        comprende all'inizio e che reggono tutto il gioco. In
        \emph{Sandcastle} i nemici non sono in scena ma possono apparire
        ovunque: horror \textbf{adrenalinico} invece che ansiogeno.
  \item Il videogioco non è lineare: creare tensione è un processo
        \textbf{iterativo} per far vivere \textbf{statisticamente} a più
        giocatori possibile lo stesso effetto; più il gioco è aperto, più è
        difficile.
  \item \textbf{Contestualizzare tutto} nelle ambientazioni (quanti bagni per
        quante persone), ma senza perdere di vista il gioco; i messaggi troppo
        complessi affidati all'ambiente spesso non arrivano.
  \item Un horror cooperativo e comico: cartoon e gruppo \textbf{riducono} la
        tensione, che torna quando si resta soli; meno regole, caos
        incentivato ma controllato; critica sociale (il postino sacrificabile);
        primo titolo \textbf{autopubblicato}.
  \item Marchio di fabbrica: derivativi ma originali, \textbf{nostalgia} degli
        anni Ottanta--Novanta, tutto fatto internamente.
  \item Le aspettative: mascherare bene i limiti fa pensare a un tripla A;
        \textbf{il prezzo deve essere coerente} con i valori produttivi. Le
        critiche utili sono quelle negative e costruttive.
  \item Obiettivi: consolidare lo studio, offrire certezza in un'industria
        instabile, crescere ``alle giuste dimensioni'' senza diventare troppo
        grandi.
\end{itemize}
